% Transcription — fonds Alexandre Grothendieck, shelfmark 156-2, batch 1 (pages 1-18).
% Established from the facsimile archives/batches/156-2/batch-01.pdf (a mirror of
% https://grothendieck.umontpellier.fr/156-2.pdf), per the skill
% transcribe-grothendieck. The folder is eighteen pages and this is its only
% batch.
%
% All eighteen leaves are mathematics and all are transcribed; there is no
% cover and no unrelated verso. His own pagination 1-18 runs across the top of
% the leaves and coincides with the archivists'. Page 1 carries the chapter
% title « Réalisations topologiques des réseaux (Juin 86) » and, in a chevron
% in the corner, « GF II » — presumably Géométrie des formes, chapter II, after
% folder 156-1. Two dated headings: « 1. Faille topologique (6 Juin) » on
% page 1 and « 7 Juin » against « 2) Boucle topologique » on page 14. The
% inventory's bracketed « [Chapitre] » is the archivists'; the title itself is
% his.
%
% Content. Pages 1-13: a closed part A of X is « régulièrement immergé » when
% the point a of X/A has a conical neighbourhood V ≅ Cône(S); the boundaries
% S_t of the tubular neighbourhoods W_t cut X into interior and exterior. A
% « faille topologique » is (X, {A,B}) with A, B closed, disjoint, regularly
% immersed, (X∖A)/B and (X∖B)/A connected (variant c': every component of
% U = X∖(A∪B) meets both A and B); the « bouts » are A and B, and the faille
% is « réduite » when they are points. A « lieu » is a closed C separating X∖C
% into a part containing A and a part containing B; the lemma of page 6
% (conical neighbourhoods glue along Z' ∩ Z'' = {c}) supplies the missing
% regularity of the composed failles F_{A,C}, F_{B,C}. Page 7: a lemma that D
% is a lieu of F_{B,C} iff C is a lieu of F_{A,D}, proved by
% X = X_{a,c} ⊔_c X_{c,d} ⊔_d X_{d,b}. Pages 8-9: disjoint lieux, oriented
% failles, the order C ≤ D on lieux (least and greatest element, filtering
% conditions b/b', divisibility c), and the equivalence between failles with
% n disjoint interior lieux and chains of n+1 oriented failles glued by
% homeomorphisms. Page 10: which compact triangulated spaces carry a faille
% structure; two claims are retracted in the margin (« Faux ! idiot ! »,
% « Faux … déjà dans le cas des segments ! »). Pages 11-13: restricting A, B
% and the lieux by a property P — a point (segments), B^n (the slab
% I × B^{n-1}), S^n — and a product-neighbourhood condition C × ]-1,1[.
% Page 12's first half, a second statement of the page-7 lemma, is struck
% through by three diagonals. Pages 14-18: « Rives » of A in X,
% Riv(A,X) = colim_V Comp(V∖A) = colim_V Γ(V, i_*F_2), injecting into
% Comp(Ã) (bijective for X paracompact); bi-rives {ρ,ρ'} with ρ + ρ' = 1;
% a « maille » with « nœud » A and bi-rive β is a structure whose cut
% X̄ = Dec_β(A,X) is a faille with ends A_ρ, A_ρ'; compatibility of two
% mailles (A_i, β_i), claimed symmetric at the top of page 18, where the text
% stops.
%
% Reading hazards fixed once here. « Riv » and « Rives » both occur and are
% kept. X̄ for the cut space and à for its preimage of A are his. His
% underlining is set in bold. The hand is drafting throughout: the formulas
% carry and the slanted marginal notes of pages 1, 3, 4 and 6 (at forty to
% sixty degrees, some hatched over in pencil) are only partly read; what they
% will not support is \ill{}.
%
% Readings flagged and not settled: « compact » for the cone base S (p. 1);
% the gloss « cl ∖ i » on page 2; « variété » in NB (3) on page 2; the
% sentence introducing the equivalence of categories on page 8, broken by
% strikings; « intrinsèque » on page 11; « ord. » closing page 11; the
% description of the boolean algebra of rives on page 15; « s'identifie » in
% a) on page 17. One variance is on the page and not repaired: page 12's
% sketch labels a lieu D^1 where B^1 is expected. Page 10's bracket (the
% structure on a segment is unique) and its margin (the family A = [0,α],
% B = [β,1]) are read as speaking of reduced and non-reduced structures
% respectively. Page 4's margin reads « mais pas néc. », not « jamais ». The lemma of page 6 does not state
% Z = Z' ∪ Z''.
%
% Pass: Opus 5 (claude-opus-5), 2026-09-13 — first pass, unchecked against the pages by a human.
\documentclass[11pt,a4paper]{article}
% Le préambule commun à toutes les transcriptions du fonds Grothendieck.
%
% Il tient en une idée : **l'incertitude se compose, elle ne se résout pas.**
% Une transcription qui lisse un mot douteux a détruit la seule chose pour
% laquelle on la faisait — un lecteur doit pouvoir savoir, à chaque endroit,
% ce qui était sur le feuillet et ce qui a été ajouté. Les six macros qui
% suivent sont donc l'appareil critique, et non de la décoration.
%
% Les mêmes commandes sont reconnues par `scripts/render.mjs`, qui produit la
% vue de lecture du site. Une macro ajoutée ici doit l'être aussi là-bas,
% faute de quoi le rendu échouera bruyamment — ce qui est voulu : un rendu
% silencieusement faux est pire qu'un rendu absent.

\ProvidesPackage{grothendieck}[2026/08/08 Transcriptions du fonds Grothendieck]

\usepackage{amsmath,amssymb,amsthm}
\usepackage{mathrsfs}
\usepackage{stmaryrd}
% Les diagrammes commutatifs sont partout dans ce fonds : c'est la langue
% même de la géométrie algébrique de Grothendieck, et une transcription qui
% les rendrait en prose ne transcrirait rien.
\usepackage{tikz-cd}
% Français par défaut — c'est la langue des feuillets — mais l'anglais est
% chargé aussi : la traduction et le résumé sont des documents anglais, et la
% typographie française (espaces avant les deux-points) y serait une faute.
% Ces documents basculent par \selectlanguage{english} en tête de corps.
\usepackage[english,french]{babel}
\usepackage{fontspec}
\usepackage{geometry}
\geometry{a4paper,margin=25mm,marginparwidth=28mm}
\usepackage{marginnote}
\usepackage{xcolor}
% ulem, not soul. `soul`'s \st and \ul are built on a character-by-character
% scan that XeLaTeX's font machinery breaks: the document fails with an
% undefined control sequence rather than degrading. ulem does the same two jobs
% and is Unicode-engine safe. `normalem` stops it hijacking \emph, which the
% transcriptions use for Grothendieck's own emphasis.
\usepackage[normalem]{ulem}
\usepackage{hyperref}
\hypersetup{colorlinks=true,linkcolor=black,urlcolor=black,pdfborder={0 0 0}}

% --- Métadonnées du lot -----------------------------------------------------
% Reprises telles quelles par le script de rendu, qui les affiche en tête de
% la vue de lecture. Elles fixent ce que le fichier prétend couvrir : sans
% elles, une transcription flotte sans point d'attache dans un fonds de seize
% mille feuillets.

\newcommand{\folder}[1]{\gdef\grothFolder{#1}}
\newcommand{\batch}[1]{\gdef\grothBatch{#1}}
\newcommand{\leaves}[2]{\gdef\grothFirst{#1}\gdef\grothLast{#2}}
\newcommand{\foldertitle}[1]{\gdef\grothFoldertitle{#1}}
\gdef\grothFolder{?}\gdef\grothBatch{?}\gdef\grothFirst{?}\gdef\grothLast{?}\gdef\grothFoldertitle{}

% --- Le résumé --------------------------------------------------------------
%
% Il ouvre la lecture modernisée, sous le titre « Résumé », et il est composé
% en petit. Ce n'est pas un ornement : c'est le seul passage du document qui
% ne soit pas des mathématiques, et un lecteur doit pouvoir le reconnaître
% avant de l'avoir lu — puis le sauter s'il connaît déjà le sujet, ou s'y
% arrêter s'il ne le connaît pas. Le filet qui le suit marque la reprise du
% corps.

\newenvironment{resume}
  {\small\setlength{\parskip}{0.45em}}
  {\par\medskip\hrule\medskip}

% --- Mots-clés ---------------------------------------------------------------
%
% En anglais, en clôture du résumé de la lecture modernisée : le vocabulaire
% moderne sous lequel chercher ce que les feuillets construisent. Le manifeste
% du site (`npm run manifest`) les extrait pour étiqueter la cote entière —
% cette ligne est la source unique des tags, il n'y a pas de fichier à côté.
\newcommand{\keywords}[1]{\par\smallskip\noindent
  {\footnotesize\textsc{Keywords} --- \textit{#1}}\par}

% --- L'appareil critique ----------------------------------------------------

% Le numéro de feuillet, en marge. C'est l'ancre qui permet de régler un
% désaccord sur une formule en regardant *un* feuillet plutôt qu'une chemise
% de six cents. Le site s'en sert aussi pour tourner les pages du fac-similé
% quand on fait défiler la transcription.
\newcommand{\leaf}[1]{%
  \marginnote{\footnotesize\color{gray!70}\textsf{#1}}%
  \label{leaf:#1}\ignorespaces}

% Illisible : on ne devine pas. Un mot inventé qui se lit comme les autres est
% le pire résultat possible.
\newcommand{\ill}{\textcolor{red!60!black}{[\,\ldots\,]}}

% Lecture douteuse : le mot est donné, et signalé comme incertain.
\newcommand{\uncertain}[1]{\uline{#1}}

% Ajout de l'éditeur — une lettre manquante, un mot sous-entendu.
\newcommand{\add}[1]{\textcolor{blue!50!black}{[#1]}}

% Biffé par Grothendieck lui-même. Ce qu'il a rayé fait partie du document :
% une rature dit souvent où la pensée a changé de direction.
\newcommand{\struck}[1]{\sout{#1}}

% Note du transcripteur, sur l'état matériel du feuillet.
\newcommand{\note}[1]{\par\smallskip{\small\color{gray!60!black}\textit{[#1]}}\par\smallskip}

% Note marginale de Grothendieck, distincte du corps du texte.
\newcommand{\marginal}[1]{\par\smallskip\noindent\hspace{1em}%
  {\small\color{gray!60!black}#1}\par\smallskip}

% --- Titre ------------------------------------------------------------------

\AtBeginDocument{%
  \begin{center}
    {\large\bfseries Fonds Alexandre Grothendieck --- cote n\textsuperscript{o}~\grothFolder}\\[2pt]
    {\itshape\grothFoldertitle}\\[4pt]
    {\small feuillets \grothFirst--\grothLast\ (lot \grothBatch)}\\[2pt]
    {\footnotesize\color{gray!60!black}Transcription établie d'après le fac-similé numérisé
     par l'Université de Montpellier.\\
     Les lectures incertaines sont soulignées, les illisibles notées [\,\ldots\,],
     les ajouts entre crochets.}
  \end{center}
  \vspace{1em}}

\endinput


\folder{156-2}
\batch{1}
\pages{1}{18}
\dating{1986}
\watermark{Édition de démonstration}
\foldertitle{[Chapitre] II. Réalisations topologiques des réseaux : notes manuscrites (06/06/1986)}

\begin{document}
\page{1}
\section{Réalisations topologiques des réseaux}

\note{titre de sa main, en tête de la page 1, suivi de « (Juin 86) » ; dans l'angle supérieur droit, encadré d'un chevron, « GF II »}

\subsection{1. Faille topologique}

\marginal{(6 Juin)}
Soit $X$ un esp. top., $A \subset X$ partie fermée \add{non vide}\note{« non vide » est écrit au-dessus de la ligne, de sa main} de $X$. \struck{\ill{}} $X/A$ l'espace déduit de $X$ en contractant $A$ en un pt , $a$ le pt déduit de $A$ par contraction\note{cette précision est écrite au-dessus de la ligne}, [Si \struck{$X$} $X'$ \struck{est} une partie de $X$ contenant $A$, alors $X'/A \hookrightarrow X/A$ identifie $X'/A$ à un sous-espace top. de $X/A$ [les fermés de $X'/A$ s'identifient aux fermés $Y'$ de $X'$ qui ou bien contiennent $A$, ou bien ne le rencontrent pas (i.e. tels que $Y' \cap A \in \{\emptyset, A\}$), et ce sont les traces sur $X'$ des fermés $Y$ de $X$ qui ont ces propriétés, lesquels correspondent d'ailleurs aux fermés de $X/A$…]] On dit que $A$ est \textbf{régulièrement immergé} dans $X$ si $\exists$ voisinage \struck{\ill{}} conique de $a$ i.e. voisinage \add{(fermé)} $V \overset{h}{\simeq} \mathrm{C\hat{o}ne}(\mathcal{S})$, $\mathcal{S}$ espace topologique \uncertain{compact}, $h(a) = o$ (sommet du cône) et $\uncertain{h}(V^{0}) = V \smallsetminus B$ ouvert dans $X/A$. \struck{Si $V$,} Alors \struck{les} sur $V$ il y a (via $h$) une structure conique, \uncertain{par laquelle} on considère les
\[
  tV = V_t, \qquad 0 \leq t \leq 1,
\]
qui sont, pour $t \in {]0,1]}$, homéom. à $V$, et forment
\marginal{\textbf{NB} $\dot V = B$, $\dot V_t = B_t$ et \ill{}}
\marginal{$\mathcal{S}$ \uncertain{connu} « section » de $V$, et \ill{}}
\note{le coin inférieur gauche porte deux notes obliques, en partie recouvertes de hachures au crayon ; la seconde n'est lue qu'à moitié. La lettre $\mathcal{S}$ est écrite en surcharge, sur ce qui semble être un $B$}

\page{2}
\marginal{$V_t = V_t^{0} \sqcup B_t$}
un système fond. de voisinages de $a$ dans $X/A$. Les $V_t^{0} = t V^{0}$, où $V^{0} = V \smallsetminus B$, forment un syst. fond. de voisinages\add{ ;} les images inverses $W_t$ de $V_t$ dans $X$ forment un syst. fond. de voisinages de $A$ dans $X$. Pour tout tel $W_t$, \struck{\ill{}} $\mathcal{S}_t = \dot{W}_t$ \marginal{cl$\smallsetminus$i}\note{petit croquis au-dessus de la ligne ; la glose « cl $\smallsetminus$ i », reliée à $\dot W_t$ par un trait, est d'une lecture incertaine} est un fermé de $X$, et \struck{\ill{}} posant $W'_t = X - V_t$, on a
\[
  X \smallsetminus \mathcal{S}_t \;=\; \underbrace{W_t^{0}}_{\text{int. de } W_t} \sqcup \underbrace{W'_t}_{\text{ext. de } W_t} \qquad \text{somme topologique}
\]
\note{sous $\mathcal{S}_t$, de sa main : « bord de $W_t$ »}

\textbf{NB} (1) L'hyp. faite sur $A \subset \struck{X}$ est \textbf{locale} i.e. \struck{Soit} si $X'$ voisinage de $A$ dans $X$, elle est vraie pour $(A,X)$ ssi elle est vraie pour $(A,X')$. (2) [Elle a un sens, et est vraie, pour $A = \emptyset$, $X/\emptyset = \{a\} \sqcup X$, on a donc $B = \emptyset$] (3) Vraie pour $A$ \add{sous-}\uncertain{variété} \ill{} \struck{\ill{}} \ill{} dans une variété $X$.
\marginal{(4) $A$ \ill{} $B = \emptyset$}

\textbf{Définition} \struck{Une} Une \textbf{faille topologique} est un couple $(X, \{A,B\})$ formé d'un espace topologique $X$, et d'une paire $\{A,B\}$ \ill{} de parties de $X$.

\page{3}
satisfaisant aux conditions suivantes

\marginal{Mais $X \neq A \cup B$ car $A$, $B$ non ouverts, donc $U = X \smallsetminus (A \cup B)$ est non vide}
\marginal{Conditions des failles réduites : a) $a$ et $b$ \struck{pts fermés} (non isolés) ; b) $a \neq b$ ; c) $X \smallsetminus \{a\}$ et $X \smallsetminus \{b\}$ connexes ; d) $a$ et $b$ ont des voisinages coniques.}
\marginal{\textbf{NB} Si $(X, \{A,B\})$ est une faille et $C \subset U = X \smallsetminus (A \cup B)$ \uncertain{fermé} \ill{}}

\begin{itemize}
  \item[Fi a)] $A$ et $B$ sont \add{fermés, non \ill{} (\uncertain{donc} non vides)} \struck{régulièrement immergés} \struck{dans} $X$. [Dans le pt $a$ de $X/A$ et le pt $b$ de $X/B$ sont non isolés, i.e. $S_A \neq \emptyset$, $S_B \neq \emptyset$]
  \item[Fi b)] \struck{\ill{}} $A \cap B = \emptyset$
  \item[Fi c)] $(X \smallsetminus A)/B$ et $(X \smallsetminus B)/A$ sont connexes
  \item[Fi d)] $A$ et $B$ sont régulièrement immergés. \struck{e) \ill{} on va formuler \ill{}}
\end{itemize}

On va donner \struck{\ill{}} une variante de c) :

\struck{ainsi} Pour

\begin{itemize}
  \item[c')] Toute composante connexe $U_i$ de $U$ ($U = X \smallsetminus (A \cup B)$) rencontre $B$ et rencontre $A$ \marginal{$\overline{U_i}$}
\end{itemize}

On a c') $\Longrightarrow$ c), et l'inverse est vrai si $U$ n'a qu'un nb fini de composantes connexes (cas $X, A, B$ « modérés », \uncertain{cas usuel}).

$A$ et $B$ sont les \textbf{bouts} \add{(extrémités)} de la faille.

La faille est \textbf{réduite} \marginal{si $A$, $B$ sont réduits à des pts $a$, $b$. On dit que $(X, \{a,b\})$ est une faille} si $A$, $B$ sont réduits \add{à} un pt. Soit

$X' = {}_{A}X/B$ déduit de $X$ en contractant $A$, $B$ en des pts, $A \mapsto a$, $B \mapsto b$, $a, b \in X'$. Alors $(X, \{A,B\})$ est une faille ssi $\{X', \struck{\ill{}} \{a,b\}\}$ est une faille.

\textbf{NB} Faille réduite est \textbf{connexe}

\page{4}
\marginal{\textbf{NB} Pour une faille réduite $(X', \{a,b\})$, \struck{\ill{}} $X'$ est connexe, ainsi que $X' \smallsetminus \{b\}$ et $X' \smallsetminus \{a\}$ (denses dans $X'$), mais pas \uncertain{néc.} $X' \smallsetminus \{a,b\}$ (ex. $S^1$ et deux pts)}
\note{sous cette note, un petit cercle marqué de deux points}

\textbf{Lieu} sur une faille $(X, \{A,B\})$ : c'est une partie \add{fermée} $C \subset X$ telle que l'on ait les propriétés : ou bien $C \in \{A,B\}$, ou bien les propriétés suivantes sont satisfaites :
\struck{Li a) $C$ fermé, et $C \in \{A,B\}$ ou $C \subset U = X \smallsetminus (A \cup B)$}

\struck{\ill{}} (\ill{} « catégorisations »…\ \ill{})
\note{deux lignes barrées d'un long trait oblique, à peine lisibles}

\begin{itemize}
  \item[Li a)] $\exists$ décomposition de $X \smallsetminus C$ \add{espace} en \uncertain{somme} de deux sous-espaces
  \[
    \mathcal{U}_{A,C}, \quad \mathcal{U}_{B,C}, \qquad A \subset \mathcal{U}_{A,C}, \quad B \subset \mathcal{U}_{B,C}
  \]
  \struck{\ill{}} et $\mathcal{U}_{A,C}/A$ et $\mathcal{U}_{B,C}/B$ connexes [et ceci détermine $\uncertain{\mathcal{U}_{A,C}}, \uncertain{\mathcal{U}_{B,C}}$ de façon unique, \uncertain{comme} les composantes connexes de $X' \smallsetminus C$ \struck{$X \smallsetminus X/B$} qui contiennent
\end{itemize}

\marginal{i.e. $X' \smallsetminus C$ a deux composantes connexes, dont l'une contient $a$, l'autre contient $b$}
\marginal{\textbf{NB} \struck{\ill{}} quand $A$, \struck{$B$} $B$ connexes, $\mathcal{U}_{A,C}$, $\mathcal{U}_{B,C}$ sont \ill{} composantes \ill{}}
\note{les deux notes précédentes sont écrites en oblique dans la marge gauche, la seconde en partie raturée}

\page{5}
\textbf{NB} Soit $X_{A,C} = \overline{\mathcal{U}_{A,C}} \cup C$, $X_{B,C} = \overline{\mathcal{U}_{B,C}} \cup C$.

Je dis que $(X_{A,C}, \{A,C\})$ satisfait les conditions a) b) c) des failles (donc il en sera de même de $(X_{B,C}, \{B,C\})$).

Pour le voir, on est ramené au cas $A, B, C$ réduits à des pts, $a, b, c$.

Il faut prouver a) \struck{\ill{}} $c$ n'est pas isolé dans $X_{a,c}$ \marginal{$a$ n'est pas isolé dans $X_{a,c}$ puisqu'il ne l'est pas dans $X$ et $X_{a,c}$ \uncertain{contient un} vois. \uncertain{de $a$})}. \struck{On} Sinon on aurait $\mathcal{U}_{ac}$ ouvert et fermé dans $X$, $X$ serait non connexe, absurde.

b) $X_{a,c} \smallsetminus \{c\}$ est connexe par hyp. Je dis que $X_{a,c} \smallsetminus \{a\}$ est connexe. Sinon, \uncertain{soit} $\exists$ partie $V$ de $X_{a,c} \smallsetminus \{a\}$ qui soit ouverte et fermée dans $X_{a,c} \smallsetminus \{a\}$, non vide, et ne contenant pas $c$. Donc elle serait aussi \struck{ouverte et} (fermée et) \textbf{ouverte} dans $X \smallsetminus \{a\}$, \struck{\ill{}} \ill{}, $X \smallsetminus \{c\}$ ne serait pas connexe, absurde.

Pour que $(X_{A,C}, \{A,C\})$ soit une faille, il ne manque que l'hypothèse d), d'immersion locale.

\page{6}
\textbf{Lemme} Soit $Z$ un espace, $Z'$ et $Z''$ deux parties fermées, telles que $Z' \cap Z'' = \{c\}$. Pour que $c$ ait un voisinage conique dans $Z$, il faut et il suffit qu'il ait des voisinages coniques dans $Z'$ et $Z''$.
\note{l'énoncé ne dit pas $Z = Z' \cup Z''$ ; c'est pourtant l'usage qui en est fait plus bas}

\ill{} seulement, pour \ill{} d'immersion de la faille, la condition \struck{\ill{}} « $C \in \{A,B\}$ ou $C$ régulièrement immergé dans $X$ ».

\textbf{Exemples} (cf. p. \ill{}) Pour $t$ assez petit, les $\dot V_{A,t} = S_{A,t}$ sont des lieux de la faille.

\textbf{Lieux disjoints.} Soit $F = (X, \{A,B\})$ une faille, $\mathcal{L}(F)$ l'ens. des lieux. Je vais définir une relation symétrique \struck{\ill{}} « disjoints » $(\Rightarrow (C \neq D))$ sur l'ens. des lieux. Si $C$ et $D$ sont \ill{} \struck{\ill{}} $C$ ou $D$ sont $\in \{A,B\}$, « disjoints » signifie simplement \textbf{distincts} ($\Longleftrightarrow$ disjoints au sens $C \cap D = \emptyset$)

\marginal{\textbf{NB} Pour qu'une partie fermée $C$ de $X$, $C \subset U = X \smallsetminus (A \cup B)$, soit un lieu, il faut et il suffit que $C \neq \emptyset$ et que \struck{donc} \ill{} $\bigl((X/A,B,C, \{a,b\}), \{c\}\bigr)$ \ill{}}
\marginal{Si $D$ est une partie fermée de $X$ telle que \struck{\ill{}} $\mathcal{U}_{A,C}$ ou $D \subset \mathcal{U}_{B,C}$ \struck{\ill{}}, alors $D$ \ill{} les \ill{} faille \ill{} lieu \ill{}}
\note{deux notes obliques remplissent le coin inférieur gauche, séparées par un trait ; plusieurs mots sont raturés en boucle et n'ont pas été lus}

\page{7}
\textbf{Décomposition} d'une faille $F = (X, \{A,B\})$ par un lieu $C$ distinct des extrémités : failles $F_{A,C}$, $F_{B,C}$. Tout lieu d'une des failles composantes est un lieu de $F$, \struck{\ill{}} (pour \textbf{homéom.})
\note{la suite de la ligne est une énumération biffée d'un seul trait, reprise en interligne :}
\struck{Équiv. de catégories \ill{}} a) Failles avec lieu intérieur, avec homéo \uncertain{b)} Paires de failles, avec un bout de l'une, un bout de l'autre. \struck{c'est \ill{}} \struck{Soient}

\textbf{Lemme} Soient $C$ et $D$ deux lieux de $F$, distincts des extrémités, avec $(C \neq D)$.
\marginal{$A \; C \; D \; B$}
\note{en marge, les quatre lettres alignées, $C$ et $D$ reliés par des crochets au-dessus et au-dessous}
Alors $D$ est un lieu de $F_{B,C}$ ssi $C$ est un lieu de $F_{A,D}$. ($D$ \uncertain{est}, \uncertain{une fois} \ill{} \uncertain{les rôles} de $A$ et $B$, \struck{$D$ \ill{}} \ill{} de $C$ et $D$, $D$ est un lieu de $F_{A,C}$ ssi $C$ est un lieu de $F_{B,D}$.)

\textbf{Démonstration} : Il suffit de prouver $\Longrightarrow$ \struck{\ill{}} (l'autre implication en résulte, en \ill{} \ill{} A et B, $C$ et $D$) \struck{\ill{}} OPS $A$, $B$, $C$, $D$ \textbf{réduits} \struck{\ill{}} à des points. On a par hyp \struck{(c lieu de $X$)}
\[
  X = X_{a,c} \amalg_c \underbrace{X_{c,b}}_{\overset{\text{déf}}{=} X_{b,c}}
\]
($c$ lieu de $X$) et
\[
  X_{c,b} = X_{c,d} \amalg_d \underbrace{X_{d,b}}_{\overset{\text{déf}}{=} X_{b,d}}
\]
(car $d$ lieu de $X_{c,b}$), donc
\[
  X = \underbrace{X_{a,c} \amalg_c X_{c,d}}_{X_{a,d}} \amalg_d X_{d,b}
\]

\page{8}
\textbf{Lieux disjoints} sur une faille : $C, D$ \marginal{(relation symétrique)}

\struck{\ill{} lieu $C$} Si $C$ ou $D$ est une extrémité, \textbf{disjoint} signifie que $C \neq D$. Si ni $C$ ni $D$ n'est une extrémité, alors disjoint signifie que \struck{\ill{}} $C \neq D$, \textbf{et} que $D$ est un lieu de $F_{A,C}$ ou de $F_{B,C}$ ($\Longleftrightarrow$ $C$ est un lieu de $F_{B,D}$ ou de $F_{A,D}$).

\textbf{Failles orientées} : Failles avec extrémité fixée (« origine »).

\struck{\ill{}} La catégorie des failles orientées avec \uncertain{deux} lieux internes disjoints, de deux \ill{} : celle des \struck{systèmes} \ill{} est équivalente \ill{} \struck{triples} $(F_1, F_2, F_3, h_{21}, h_{32})$, où $F_1, F_2, F_3$ sont des failles orientées,
\[
  h_{21} : \mathrm{ex}(F_1) \simeq \mathrm{or}(F_2), \qquad h_{32} : \mathrm{ex}(F_2) \simeq \mathrm{or}(F_3).
\]
\note{« ex » et « or » : extrémité et origine d'une faille orientée ; l'énoncé est interrompu par des biffures et sa première ligne n'est lue qu'en partie}

Soit $F = (X, (A,B))$ faille orientée. Soit $C$ un lieu, on pose $\mathcal{F}_C =$ ens. des lieux de $F_{A,C}$
\marginal{si $C \neq A, B$ ; $= \{A\}$ si $C = A$ ; $= \ill{}$ si $C = B$}.
On écrit $C \leq D$ ssi $\mathcal{F}_C \subset \mathcal{F}_D$. Cela signifie simplement que $C \in \mathcal{F}_D$.

\page{9}
\marginal{Alors \ill{} $F_{A,C} \subset F_{A,D}$ \ill{} $F_{A,C} = F_{A,D} \Rightarrow C = D$ \ill{}}
Les lieux d'une faille \add{orientée} forment un ens. $\mathcal{L}$ ordonné. \uncertain{Il satisfait} les conditions suivantes :
\begin{itemize}
  \item[a)] $\exists$ plus petit et plus grand élément $\alpha$ et $\beta$, distincts
  \item[b)] $\forall x \in \mathcal{L}$, $x \neq \beta$, \struck{\ill{}} l'ens. des élts $> x$ est filtrant décroissant ;
  \item[b')] $\forall x \in \mathcal{L}$, $\alpha \neq x$, l'ens. des élts $< x$ est filtrant croissant
  \item[c)] Si $x < y$, $\exists z$ avec $x < z < y$ (divisibilité)
\end{itemize}

\textbf{NB} (1) Deux lieux $x, y$ \uncertain{sont} \textbf{disjoints}, \struck{ssi} ssi $x < y$ ou $y < x$. Ils définissent \uncertain{$F_{x,y} = F_{y,x}$}, \uncertain{alors} \ill{} une faille \ill{}, dont les lieux intermédiaires, naturellement \ill{} (\ill{} forment une partie \textbf{totalement} ordonnée).
\marginal{(2) Si on change l'orientation, on remplace l'ordre par l'ordre opposé}

La catégorie \ill{} d'une faille orientée $F$, et de $n$ tels lieux $x_1, \dots, x_n$ \uncertain{disjoints}, est équivalente : celle des systèmes $(F_1, \dots, F_{n+1}, h_{21}, h_{32}, \dots, h_{n+1,n})$ de $n+1$ failles orientées, et de $n$ homéomorphismes de recollement.

\page{10}
\subsection{Conditions particulières sur les lieux.}

\textbf{NB} Pour qu'un espace compact triangulé admette une structure de faille, il f. et s. qu'il soit connexe, et non réduit à un point. C'est gros ! Il y a alors $\infty$ \uncertain{de} structures de faille réduite. [NB. Si $X$ est un segment, \ill{} cette str. est unique. C'est le seul cas \ill{} !]
\marginal{\textbf{NB} Les structures de faille sur $[0,1]$ \ill{} $\{A,B\}$ \ill{} données par $A = [0,\alpha]$, $B = [\beta,1]$, $0 \leq \alpha < \beta \leq 1$}
\note{la note marginale décrit les structures de faille non réduites du segment ; l'unicité du crochet porte, semble-t-il, sur les structures réduites}

On peut considérer rel. d'équivalence entre structures de faille \struck{sur $X$} (resp. failles orientées) sur $X$ : \ill{} celle engendrée par la relation d'ordre (\uncertain{cas orienté}) \struck{$(A \subset A'$ ou $A \subset B')$ et $(B \subset A', B \subset B')$} \struck{\ill{}}
\[
  A \subset A' \quad \text{et} \quad B \subset B' .
\]
\struck{[\textbf{NB} Si $(X, \{A,B\})$ est une faille, $A'$, $B'$ sont des fermés, et si \uncertain{$A'$, $B'$} régulièrement immergés tels que $A \subset A'$, $B \subset B'$, $A' \cap B' = \emptyset$, alors $(X, \{A',B'\})$ est une faille.]}
\marginal{Faux ! idiot !}
\note{ce crochet est barré de deux longues diagonales}
\marginal{Il est \ill{} les voisinages \ill{} de $A$, $B$ \ill{} voisins \ill{}}

\struck{\ill{}} \ill{} que deux structures de faille orientée sur un $X$ triangulé sont toujours équivalentes.
\marginal{Faux (cas orienté !) déjà dans le cas des segments !}

\page{11}
\struck{Considéra\ill{}} On peut imposer des conditions à $A$, $B$, comme d'être connexes, voire contractiles, voire $\simeq$, homéomorphes à un espace donné, et itou pour les lieux, \add{d'où} [\ill{} $P$ propriété \uncertain{intrinsèque}] \ill{} \struck{\ill{}} $\mathcal{L}_P$. Pour que les conditions plus haut sur l'ensemble ordonné $\mathcal{L}$ \struck{soient} restent vraies pour $\mathcal{L}_P$, il est bon de \struck{\ill{}} (pour $P$ \uncertain{choisie} p. ex. \struck{d'\ill{}} \ill{} d'être contractile, ou homéomorphe à un espace fixé) que les espaces $S_A$, $S_B$ (bords des vois. tub. de $A$, $B$ dans $X$) aient cette propriété, et \struck{itou} on \uncertain{fixe} la restriction correspondante sur les lieux internes $C$ : les bords des vois. tub. de $C$ dans $X_{A,C}$ et $X_{B,C}$ ont la propriété $P$.

\textbf{Exemples}

1) $P$ : être \struck{homéom.} réduit à un pt. On trouve la théorie des segments top. \uncertain{ord.}

\page{12}
\note{toute la première moitié de la page (le lemme et le début de sa démonstration) est barrée de trois longues diagonales ; elle reprend le lemme de la page 7}
\struck{\textbf{Lemme} Soit $F = (X, \{A,B\})$ une faille, $C$ un lieu de \ill{} $F$ distinct des extrémités, \ill{} des \ill{} sous-divisions $F_{A,C} = (X_{A,C}, \{A,C\})$, $F_{B,C} = (X_{B,C}, \{B,C\})$. Soit $D$ un \ill{} lieu de $F_{B,C}$, distinct des extrémités. Alors $C$ est un lieu de $F_{A,D}$, distinct des extrémités, et inversement.}

\marginal{$A \; C \; D \; B$}
\note{en marge, un croquis : une bande allongée, $A$ à gauche, deux traits transverses $C$ et $D$, $B$ à droite}

\struck{OPS $A, B, C, D$ réduits à des points.} \struck{Failles orientées} \struck{$A, C, D \Longrightarrow$} \ill{}
\[
  \struck{D, B, C} \Longrightarrow \struck{C, A, D} \Longrightarrow \struck{D, B, C}
\]
\note{ces lignes sont sous les mêmes diagonales que le lemme}

\struck{1)} 2) Être homéom. à $B^n$ : On trouve notamment la théorie de la décomposition de
\[
  I \times B^{n-1} \simeq B^n
\]
(avec les deux « extrémités » \struck{$B^{n-1} \times \{0\}$}, $\{0\} \times B^{n-1}$, $\{1\} \times B^{n-1}$) par des « lieux » homéomorphes à $B^{n-1}$.

\note{croquis à gauche du texte : un rectangle hachuré, ses deux côtés verticaux épaissis et marqués $A \simeq B^1$ et $B \simeq B^1$, traversé d'une courbe marquée $C \simeq D^1$ ; au-dessous, « ou », un fuseau hachuré entre deux points $A$ et $B$, traversé de même. $D^1$ est écrit ainsi, là où l'on attend $B^1$}

\page{13}
3) Être homéom. à $S^n$. On trouve p. ex. sur $S^{n+1}$ muni des deux pts \add{distincts} (p. ex. antipodiques, si on veut) $a$, $b$, une théorie de la décomposition par des $S^n$.

\note{croquis : un cercle portant deux points marqués $a$ et $b$, et coupé en haut et en bas par deux petits traits marqués chacun $c$ — le lieu est ici une sphère $S^0$, deux points}

On peut aussi exiger, pour un lieu distinct des extrémités, que la structure de $X$ au voisinage de $C$ soit celle de $C \times {]-1,1[}$, avec $C$ correspondant à $C_0 = C \times \{0\}$, et les deux « rives » $C \times {]-1,0[}$, $C \times {]0,1[}$ étant contenues resp. dans les deux failles composantes $X_{A,C}$, $X_{B,C}$.

\page{14}
\subsection{2) Boucle topologique}

\note{le titre porte deux traits, dont le second pourrait être une rature ; le paragraphe qui suit ouvre sur un autre titre, souligné}

\marginal{7 Juin}
\struck{Soit} \textbf{Rives d'une partie fermée} $A$ d'un espace $X$. \struck{Considérons les} Soit
\[
  Z \mapsto \mathrm{Comp}(Z) = \mathrm{Hom}_{\mathrm{cont}}(Z, \{0,1\}) \;\bigl(= \Gamma(Z, \mathbb{F}_{2,Z})\bigr)
\]
(l'ens. des parties \uncertain{ouvertes et fermées}) de $Z$. On fait un changement \ill{} sur $Z$. On pose
\[
  \underbrace{\mathrm{Rives}(A,X)}_{\text{ens. des rives de } A \text{ dans } X} = \varinjlim_{V \text{ voisinage de } A \text{ dans } X} \mathrm{Comp}(V \smallsetminus A) = \varinjlim_{V} \Gamma(V, i_*(\mathbb{F}_{2,U}))
\]
\marginal{où $i : U = X \smallsetminus A \hookrightarrow X$ l'inclusion}

\textbf{NB} 1) Supposons que $A$ soit régulièrement immergé : prenons $S_{A,X}$ voisinage tubulaire. Alors on a
\[
  \mathrm{Riv}(A,X) \simeq \mathrm{Comp}(S_{A,X})
\]
\note{$S_{A,X}$ désigne, comme aux pages 1 et 11, le bord d'un voisinage tubulaire de $A$}

2) Supposons que $X$ soit paracompact, $F$ un faisceau sur $X$, \uncertain{alors} pour \uncertain{$A$ fermé} on a
\[
  \Gamma(A, F|A) = \varinjlim_{V \in \mathcal{V}(A,X)} F(V)
\]

\struck{Alors on a} Considérons le découpage de $X$ suivant $A$, $\widetilde{A} = \mathrm{Spec}(i_*(\mathbb{F}_{2,U}))$, puis $(\widetilde{X}, \widetilde{A})$, où $i : U = X \smallsetminus A \hookrightarrow X$. On a
\[
  \mathrm{Comp}(\widetilde{A}) \simeq \Gamma(A, i_*(\mathbb{F}_{2,U}))
\]
\struck{Donc} \ill{}
\note{la page s'arrête sur deux mots raturés}

\page{15}
\uncertain{d'où} application canonique
\[
  (*) \qquad \mathrm{Riv}(A,X) \longrightarrow \mathrm{Comp}(\widetilde{A})
\]
qui est \add{injective, et} \textbf{bijective} si $X$ paracompact. Donc dans ce cas, choisir une rive, c'est choisir une composante de $\widetilde{A}$.
\note{« injective, et » est écrit au-dessus de « bijective », de sa main}

\textbf{NB} L'ens. des \struck{composantes} rives de $A$ dans $X$ \struck{ouvertes} est une $\mathbb{F}_2$-algèbre \add{idempotente} ($E^2 = E$ $\forall E \in \mathrm{Riv}$) avec unité, l'\uncertain{algèbre} \ill{} \ill{} \ill{} de $A$ dans $X$.
\marginal{\textbf{NB} Cette algèbre \ill{} $\mathrm{Riv}(A)$ \ill{} faisceau \ill{}}

Une \textbf{bi-rive} est une \textbf{paire} $\{\rho, \rho'\}$ d'éléments de $\mathrm{Riv}(A)$, telle que
\[
  \rho + \rho' = 1
\]
(\uncertain{idempotents} « décompositions en 2 morceaux de $X - A$ au voisinage de $A$ » \ill{})
\struck{c'est dire} (ce qui implique, si $A \neq \emptyset$, que $\rho \neq \rho'$). La bi-rive est triviale si $\{\rho, \rho'\} = \{0, 1\}$. On dit que $A$ disconnecte $X$ loc. si $\exists$ bi-rive non triviale.

\struck{Soit \ill{}} Une bi-rive définit $\{\widetilde{A}_1, \widetilde{A}_2\}$ telle que $\widetilde{A} = \widetilde{A}_1 \amalg \widetilde{A}_2$.

On dit que \struck{c'est une} $\{A, \struck{\{\rho,\rho'\}}\}$
\note{la page s'arrête sur cette phrase inachevée}

\page{16}
définit sur $X$ une structure de \struck{bouclage régulier} \textbf{maille}, \struck{dont $A$} \struck{et de} \textbf{nœud} $A$, si \struck{$(\widetilde{X}, \{\widetilde{A}_1, \widetilde{A}_2\})$} \struck{et} $A$ est régulièrement immergé dans $X$ (alors \struck{\ill{}} $\widetilde{A}$ est régulièrement immergé dans $\widetilde{X}$, donc $\widetilde{A}_1$, $\widetilde{A}_2$ aussi) et si $(\widetilde{X}, \{\widetilde{A}_1, \widetilde{A}_2\})$ satisfait les conditions a) b) c) des failles top. [alors c'est une faille
\marginal{réciproque}
top. si $\widetilde{A}$ est rég. immergé dans $\widetilde{X}$].
\note{« réciproque » est écrit en marge, devant un trait vertical qui embrasse les lignes de la parenthèse}

\struck{Considérons} On trouve
\[
  \widetilde{X} \amalg_{\widetilde{A}_1} (\widetilde{A}_1 \to A) \amalg_{\widetilde{A}_2} (\widetilde{A}_2 \to A),
\]
un espace \ill{} \struck{Avec} deux parties
\marginal{$\beta = \{\rho, \rho'\}$ la bi-rive}
\[
  \overline{X} = \mathrm{Dec}_{\beta}(A, X),
\]
(le découpage par $\beta$), can. isom. à $A$, \struck{\ill{}} i.e. une partie isom. à $A \times \beta$, $\{A_\rho, A_{\rho'}\}$,
\[
  \overline{X} \longrightarrow X
\]
induisant un isom
\[
  \overline{X} \smallsetminus \underbrace{\{A_\rho \cup A_{\rho'}\}}_{A \times \beta} \longrightarrow X \smallsetminus A
\]
et des homéom.
\[
  A_\rho \xrightarrow{\ \simeq\ } A, \qquad A_{\rho'} \xrightarrow{\ \simeq\ } A,
\]
i.e. la projection can.
\[
  A \times \beta \longrightarrow A
\]
\note{la page s'arrête ici ; la phrase continue à la page suivante}

\page{17}
Ceci dit, on \struck{dit} que $(\struck{A}, \beta)$ définit \add{une} une structure de maille à \textbf{nœud} $A$, à \struck{\ill{}} bi-rive $\beta$, \struck{si} si $(\overline{X}, \{A_\rho, A_{\rho'}\})$ est une \textbf{faille} top.

Orienter la maille, c'est par définition orienter la faille, i.e. choisir $A_\rho$ ou $A_{\rho'}$, i.e. choisir $\rho$ ou $\rho'$, i.e. une des rives que contient la bi-rive.
\marginal{\textbf{NB} Dans la définition, \ill{} $A_1$, $A_2$ réduits à un point…}

\marginal{$A \; \rho \; \rho' \; U \; V$}
\note{croquis à gauche : une couronne entre deux cercles, coupée en haut par un trait épais marqué $A$, les deux bords du trait marqués $\rho$ et $\rho'$, et en bas par un second trait ; les deux moitiés de la couronne sont marquées $U$ et $V$}

\struck{On dit que deux structures} \struck{de} Considérons deux \struck{\ill{}} structures de mailles $(A_1, \underbrace{\{\rho_1, \rho'_1\}}_{\beta_1})$ et $(A_2, \underbrace{\{\rho_2, \rho'_2\}}_{\beta_2})$ sur $X$. Je dis qu'elles sont \struck{\ill{}} \struck{compatibles} \textbf{compatibles} \struck{\ill{}} si
\begin{itemize}
  \item[\struck{a)}] \struck{$A \cap B$}
  \item[a)] $A_1 \cap A_2 = \emptyset$ donc $A_2$ \struck{définit} \uncertain{s'identifie} à une partie de $\overline{X}^{(A_1,\beta_1)} - A_1 \times \beta_1$,
  \item[\struck{b)}] \struck{$A_2$ si}
  \item[b)] $A_2$ est un \textbf{lieu} de \uncertain{cette faille}
  \item[c)] Si $U$, $V$ sont les deux composantes connexes de $\overline{X}^{(A_1,\beta_1)} \smallsetminus \bigl((A_1, \rho_1) \cup (A_1, \rho'_1)\bigr) \smallsetminus A_2$,
\end{itemize}
\note{la page s'arrête au milieu de la condition c)}

\page{18}
alors $\{U, V\}$ définit à la fois \struck{\ill{}} $\beta_1$, et $\beta_2$. Je dis que cette condition est en fait symétrique en les deux nœuds $(A_i, \beta_i)$.
\note{le texte s'arrête ici ; le reste de la page est blanc}

\end{document}
